\documentclass[12pt, a4paper]{article}

% --------------------------------------------------------
% 宏包引入
% --------------------------------------------------------
\usepackage[UTF8]{ctex}     % 中文支持
\usepackage{amsmath, amssymb, amsthm} % 数学公式与符号
\usepackage[margin=2.5cm]{geometry}   % 页面边距
\usepackage{xcolor}         % 颜色支持
\usepackage{fancyhdr}       % 页眉页脚
\usepackage{tcolorbox}      %漂亮的文本框
\tcbuselibrary{skins, breakable}
\usepackage{enumitem}       % 列表定制
\usepackage{mathrsfs}       % 花体字

% --------------------------------------------------------
% 样式设置
% --------------------------------------------------------
% 定义颜色
\definecolor{mainblue}{RGB}{0, 70, 140}
\definecolor{accentred}{RGB}{200, 40, 40}
\definecolor{lightgray}{RGB}{245, 245, 245}

% 页眉设置
\pagestyle{fancy}
\fancyhf{}
\lhead{\textbf{高等数学专题讲义}}
\rhead{\textit{幂指函数极限与 $e$ 指数化}}
\cfoot{\thepage}
\renewcommand{\headrulewidth}{0.4pt}

% 定义定理/公式框
\newtcolorbox{conceptbox}[1]{
  colback=blue!5!white,
  colframe=mainblue,
  fonttitle=\bfseries,
  title={#1},
  arc=0mm,
  boxrule=1pt,
  left=5pt, right=5pt, top=5pt, bottom=5pt
}

\newtcolorbox{alertbox}[1]{
  colback=red!5!white,
  colframe=accentred,
  fonttitle=\bfseries,
  title={#1},
  arc=0mm,
  boxrule=1pt
}

% --------------------------------------------------------
% 正文开始
% --------------------------------------------------------
\begin{document}

\section*{\centering \huge 专题：幂指函数极限的 $e$ 指数化处理}
\begin{center}
    \large \textbf{Calculus Advanced Lecture Notes}
\end{center}
\vspace{1em}

\section{核心知识体系 (Core Concepts)}

\subsection{1. 幂指函数的定义与痛点}
形式为 $y = u(x)^{v(x)}$ 的函数称为\textbf{幂指函数}。
在求极限 $\lim_{x \to x_0} u(x)^{v(x)}$ 时，不能直接使用幂函数或指数函数的法则。
\textbf{常见不定式类型}：$1^\infty, \quad 0^0, \quad \infty^0$。

\subsection{2. 通用解法：$e$ 指数恒等变形}
这是解决幂指函数最本质、最通用的方法，能将“底数”与“指数”分离，转化为乘积形式。

\begin{conceptbox}{核心公式：$e$ 指数化 (The $e$-Identity)}
对于任意 $u(x) > 0$，恒有：
\[
u(x)^{v(x)} = e^{\ln \left( u(x)^{v(x)} \right)} = e^{v(x) \cdot \ln u(x)}
\]
利用指数函数的连续性，极限运算可深入到指数部分：
\[
\lim u(x)^{v(x)} = e^{\lim \left[ v(x) \cdot \ln u(x) \right]}
\]
\end{conceptbox}

\subsection{3. $1^\infty$ 型的各种处理路径}
对于 $\lim u(x) = 1$ 且 $\lim v(x) = \infty$ 的情形（即 $1^\infty$ 型），有以下两种主流处理方式：

\begin{itemize}
    \item \textbf{路径 A（推荐）：直接 $e$ 指数化 + 洛必达/等价无穷小}
    \[
    \lim v \ln u = \lim v \ln[1 + (u-1)] \sim \lim v(u-1)
    \]
    此路径通用性强，不仅限于 $1^\infty$，也适用于 $0^0$ 等。
    
    \item \textbf{路径 B：利用重要极限凑形}
    \[
    \lim_{x \to x_0} u^v = e^{\lim_{x \to x_0} (u-1)v}
    \]
    \textit{注：路径 B 实质上是路径 A 在 $1^\infty$ 情形下的推论。}
\end{itemize}

\newpage

\section{典型例题精讲 (Typical Examples)}

\begin{conceptbox}{解题三步走策略}
\begin{enumerate}
    \item \textbf{定型}：判断极限类型 ($1^\infty, 0^0, \infty^0$)。
    \item \textbf{变形}：使用 $e^{v \ln u}$ 将底数化为 $e$，将问题转化为求指数部分 $A = \lim v \ln u$ 的极限。
    \item \textbf{求值}：对 $A$ 使用\textbf{洛必达法则}或\textbf{等价无穷小替换}求解。
\end{enumerate}
\end{conceptbox}

% 例题 1
\subsection*{题型一：$1^\infty$ 型与洛必达法则的结合}
\textbf{例题 1.} 计算极限 $\displaystyle \lim_{x \to 0} (\cos x)^{\frac{1}{x^2}}$。

\noindent\textbf{【解析】}
此为 $1^\infty$ 型不定式。
\begin{align*}
    \text{原式} &= \lim_{x \to 0} e^{\frac{1}{x^2} \ln(\cos x)} & \text{($e$ 指数化)} \\
    &= e^{\lim_{x \to 0} \frac{\ln(\cos x)}{x^2}} & \text{(指数连续性)}
\end{align*}
令指数部分极限为 $A$，这是 $0/0$ 型，使用洛必达法则：
\begin{align*}
    A &= \lim_{x \to 0} \frac{\ln(\cos x)}{x^2} \quad \left(\frac{0}{0}\right) \\
    &\overset{\text{L'H}}{=} \lim_{x \to 0} \frac{\frac{1}{\cos x} \cdot (-\sin x)}{2x} \\
    &= \lim_{x \to 0} \frac{-\tan x}{2x} \\
    &= -\frac{1}{2} \cdot \lim_{x \to 0} \frac{\tan x}{x} = -\frac{1}{2} \cdot 1 = -\frac{1}{2}
\end{align*}
\textbf{【结果】} 原极限 $= e^{-\frac{1}{2}}$。

\vspace{1em}
\noindent\textbf{【技巧点拨】} 也可以使用等价无穷小 $\ln(\cos x) \sim \cos x - 1 \sim -\frac{1}{2}x^2$，计算更快捷。

% 例题 2
\vspace{1.5em}
\subsection*{题型二：$\infty$ 处的幂指函数}
\textbf{例题 2.} 计算极限 $\displaystyle \lim_{x \to \infty} \left( \frac{x^2+1}{x^2-2} \right)^{x^2}$。

\noindent\textbf{【解析】}
底数 $\to 1$，指数 $\to \infty$，属 $1^\infty$ 型。
\begin{align*}
    \text{原式} &= e^{\lim_{x \to \infty} x^2 \ln \left( \frac{x^2+1}{x^2-2} \right)} \\
    &= e^{\lim_{x \to \infty} x^2 \ln \left( 1 + \frac{3}{x^2-2} \right)} 
\end{align*}
利用等价无穷小 $\ln(1+t) \sim t$ (当 $t \to 0$)，此处 $t = \frac{3}{x^2-2} \to 0$：
\[
    \text{指数极限} = \lim_{x \to \infty} x^2 \cdot \frac{3}{x^2-2} = \lim_{x \to \infty} \frac{3x^2}{x^2-2} = 3
\]
\textbf{【结果】} 原极限 $= e^3$。

\newpage

% 例题 3
\subsection*{题型三：非 $1^\infty$ 型 ($0^0, \infty^0$) 的处理}
\textbf{例题 3.} 计算极限 $\displaystyle \lim_{x \to 0^+} x^{\sin x}$。

\noindent\textbf{【解析】}
此为 $0^0$ 型不定式，\textbf{只能}使用 $e$ 指数化法，不可套用 $e^{(u-1)v}$ 公式。
\[
    \text{原式} = \lim_{x \to 0^+} e^{\sin x \cdot \ln x}
\]
考察指数部分 $A = \lim_{x \to 0^+} \sin x \cdot \ln x$ ($0 \cdot \infty$ 型)：
\begin{align*}
    A &= \lim_{x \to 0^+} \frac{\ln x}{\frac{1}{\sin x}} \quad \left(\frac{\infty}{\infty}\right) \\
    &\overset{\text{L'H}}{=} \lim_{x \to 0^+} \frac{\frac{1}{x}}{-\csc x \cot x} \\
    &= \lim_{x \to 0^+} \frac{\sin^2 x}{-x \cos x} \quad (\text{化简整理}) \\
    &= \lim_{x \to 0^+} \frac{x^2}{-x \cdot 1} \quad (\sin x \sim x, \cos x \to 1) \\
    &= \lim_{x \to 0^+} (-x) = 0
\end{align*}
\textbf{【结果】} 原极限 $= e^0 = 1$。

\section{易错警示 (Warnings)}

\begin{alertbox}{警示：乱用 $e^{(u-1)v}$ 公式}
很多同学看到 $u^v$ 就直接写成 $e^{(u-1)v}$，这是错误的！
\begin{itemize}
    \item \textbf{适用前提}：必须是 $1^\infty$ 型，即 $\lim u(x) = 1$ 且 $\lim v(x) = \infty$。
    \item \textbf{反例}：求 $\lim_{x \to 0} (1+x)^2$。
    \begin{itemize}
        \item \textbf{错误做法}：$e^{\lim 2(1+x-1)} = e^{\lim 2x} = e^0 = 1$。
        \item \textbf{正确做法}：直接代入 $(1+0)^2 = 1$。虽然此例结果巧合相同，但逻辑完全错误。如果是 $\lim_{x \to 0} 2^x \neq e^{(2-1)x}$。
    \end{itemize}
    \item \textbf{建议}：养成写 $e^{v \ln u}$ 的习惯，这是万能的。
\end{itemize}
\end{alertbox}

\section{强化练习 (Exercises)}

请利用 $e$ 指数化方法与洛必达法则计算下列极限：
\begin{enumerate}
    \item $\displaystyle \lim_{x \to 0} \left( \frac{\sin x}{x} \right)^{\frac{1}{x^2}}$
    \item $\displaystyle \lim_{x \to \infty} \left( \sin \frac{1}{x} + \cos \frac{1}{x} \right)^x$
    \item $\displaystyle \lim_{x \to 1} x^{\frac{1}{1-x}}$
\end{enumerate}

\vspace{1em}
\noindent\textbf{参考答案：} 
1. $e^{-\frac{1}{6}}$ \quad 2. $e$ \quad 3. $e^{-1}$

\end{document}